\section{表示粒度、Bidirectionality 与多轮缓存}

前一章分析参数分离，本章处理剩余两项结构：decoder 只读取 encoder 最终层，以及 encoder 使用 bidirectional attention。两项都曾提高任务质量，但随着网络更深、对话更长，信息粒度和重复计算问题会变得突出。这里第一次把 architecture choice 与 serving cost 放到同一张图中。

\subsection{只读最终 encoder layer 会不会形成信息瓶颈}

本节关注 cross-attention 的读取位置。深层网络往往在不同层表示不同粒度的信息；若 decoder 的第一层也只能访问 encoder 最终层，低层形式信息可能经过过度压缩。当前几十层模型上影响也许很小，但极深网络可能放大这种不匹配。

\lecturefigure{hyung-slide-059.jpg}{第二项假设：目标只需读取已经完全编码的输入}{Hyung Won Chung official deck, p.~59}

\lecturefigure{hyung-slide-060.jpg}{标准 cross-attention 只读取 encoder 最终层激活}{Hyung Won Chung official deck, p.~60}

\lecturefigure{hyung-slide-061.jpg}{深层网络的底层与顶层通常编码不同粒度的信息}{Hyung Won Chung official deck, p.~61}

\begin{knowledgebox}{读图：representation hierarchy 是启发式，不是固定语义表}
视觉网络常用“底层边缘、顶层物体”说明层级表示；语言模型中也可能出现局部形式到抽象语义的趋势。但层含义会随训练目标、残差连接与深度变化。图真正支持的是：不同层可能提供互补信息，因此只暴露最终层是一项需要检验的结构选择。
\end{knowledgebox}

\teachervoice{01:02:43--01:03:25，Hyung Won 说在自己用过的 24 层左右 T5 上，layer 1 读取 layer 24 似乎没问题；但若未来有十倍或千倍深度，他“不太舒服”。这是 scale extrapolation 的典型提问，而不是已观察到的失败。}

\subsection{Layer-aligned connection 是怎样的反事实}

上一节指出粒度问题，本节构造最小替代方案：让 decoder 第 $\ell$ 层读取 encoder 第 $\ell$ 层。这样每个阶段得到相近深度的输入表示，也更接近 decoder-only 同层 token interaction。反事实的价值在于可做受控实验，即使最终不采用它。

\lecturefigure{hyung-slide-062.jpg}{第三项假设：输入内部需要 all-to-all bidirectional interaction}{Hyung Won Chung official deck, p.~62}

\lecturefigure{hyung-slide-063.jpg}{Layer-aligned connection 与 causal input 的替代架构}{Hyung Won Chung official deck, p.~63}

\begin{importantbox}{结构分析的正确产物是 ablation axis}
从 H059--H063 可以得到两个独立实验轴：final-layer versus layer-aligned cross-attention，以及 bidirectional versus causal input mask。若一次同时改变两者，性能差异无法归因；课程的四步推导正是为了生成可单独缩放的 ablation。
\end{importantbox}

\subsection{Bidirectionality 的质量收益与系统成本}

Bidirectional encoder 在 BERT 时代对困难理解任务非常重要，因为每个 token 能同时利用左右上下文。但在大规模 instruction tuning 中，讲者的 anecdotal experience 是双向与单向差异不大；与此同时，多轮聊天暴露了明显系统成本：新消息改变了旧输入的“未来”，双向表示必须重新编码。

\lecturefigure{hyung-slide-064.jpg}{在足够规模下 bidirectionality 可能收益变小，却增加多轮对话成本}{Hyung Won Chung official deck, p.~64}

\teachervoice{01:03:35--01:04:23，Hyung Won 明确把结论标成 highly anecdotal：Flan-2 中双向和单向 fine-tuning 差异不大。讲义保留这个证据等级，不能把它改写成“bidirectional attention 已被证明无用”。}

\subsection{为什么 causal attention 可以复用 KV cache}

本节把图转成计算账。Causal model 中旧 token 的表示不依赖未来新 token，因此它们在前一轮计算出的 key/value 可以缓存；追加消息时只计算新增部分。Bidirectional encoder 中旧 token 可以读取新 token，追加后旧表示失效，通常需要重新编码整个输入。

\lecturefigure{hyung-slide-066.jpg}{多轮对话中 bidirectional 重编码与 unidirectional cache reuse 的对比}{Hyung Won Chung official deck, p.~66}

设已有上下文长度为 $n$，新增长度为 $m$。忽略常数与 MLP，重新做 full attention 的代价近似
\[
O\bigl((n+m)^2\bigr),
\]
而 causal KV cache（key-value cache，保存旧 token 每层的 key/value 张量）只需为新增 query 计算对旧缓存和新增 token 的 attention，增量代价近似
\[
O(nm+m^2).
\]
其中 $n$ 是旧上下文长度，$m$ 是新增 token 数。缓存降低重复计算，却增加显存占用；长上下文服务仍需管理 cache placement、batching 与 eviction。

\begin{knowledgebox}{读图：蓝色可复用区域对应因果不变性}
Unidirectional 情况下，旧 token 不能看未来，因此新消息到来不会改变旧 key/value；bidirectional 情况下，``How are you?'' 的表示会因为后续 ``Bad'' 出现而改变。系统效率优势来自计算图的依赖方向，而不是 decoder-only 名称本身。
\end{knowledgebox}

\teachervoice{01:04:31--01:05:36 的课堂解释把模型质量与 serving 连接起来：一个在 2018 年提高 benchmark 的结构，可能在多轮产品形态中制造重编码成本。Architecture trajectory 不只由 accuracy 决定，也由交互模式决定。}

\subsection{结论与 Q\&A：下一处结构瓶颈可能是 objective}

架构分析结束后，讲者把方法推回一般研究。历史结构不是无关旧物；反复做这种分析，会训练研究者识别当前问题中的隐含假设，并追问能否用更一般的方法配合规模替代它。Q\&A 进一步提出，当前更强的瓶颈也许不在 Transformer block，而在学习目标。

\lecturefigure{hyung-slide-067.jpg}{结论：识别主导力量，并用 scaling 视角分析额外结构}{Hyung Won Chung official deck, p.~67}

\teachervoice{01:06:06--01:06:42 的结尾要求学生回到自己的问题，问“哪些假设需要复查、为什么、能否换成更一般的机制并 scale up”。这比记住 decoder-only 胜出更重要，因为同一方法还可用于 MoE、multimodality、memory 与 learning objective。}

最大似然训练（maximum-likelihood estimation, MLE）把观察到的目标当作训练样本中的正确 continuation：
\[
\theta^{\star}=\arg\max_{\theta}\sum_{(x,y)\in\mathcal{D}}\log p_{\theta}(y\mid x).
\]
其中 $\mathcal{D}$ 是输入目标数据集，$x$ 是条件，$y$ 是观测目标，$\theta$ 是模型参数。对于翻译或分类，一个参考答案常能提供有效监督；对于“写一首诗”等开放任务，把单个 $y$ 当成唯一正向信号会造成很强结构假设。

\begin{knowledgebox}{RLHF 在这里扮演什么角色}
RLHF（reinforcement learning from human feedback）先从人类偏好训练 reward model，再用该奖励优化策略。Hyung Won 把它视为“学习 objective function”的一个例子：不再只复制单个目标，而让模型学习多个回答的相对质量。他同时强调 RLHF 本身也不够 scalable，因此这只是方向证明，不是终局方案。
\end{knowledgebox}

\teachervoice{01:12:39--01:14:34，Hyung Won 说自己不认为 architecture 是当前主要 bottleneck；他更不安于 MLE 对开放任务施加“只有这个 target 正确”的教学信号。这个 Q\&A 扩展了整堂课：需要被移除的 inductive bias 也可能藏在 objective 中。}

\begin{warningbox}{Compute trajectory 不是无条件承诺}
01:15:17 之后，讲者把 Moore's law 称为 red herring，强调真正变量是 compute availability；他列出 GPU、低精度与硬件专用化，也承认能源和物理可能成为瓶颈。研究计划应做多情景分析，而不是把过去指数趋势机械外推成未来保证。
\end{warningbox}

\subsection{本章小结}

表示层级提醒我们检查 final-layer bottleneck；多轮对话提醒我们把 attention direction 与缓存成本一起考虑；Q\&A 又把结构分析推广到 MLE 与 RLHF。共同方法仍是：找出当前系统中被默认接受的假设，说明它在旧 regime 中为何有用，再用规模、数据和交互变化测试它是否仍然值得保留。

